\emph{Selection-rule consistency of stacked diagrams at coincident caps.} Under the hypotheses of the selection rule transplanted in the first of the four preceding paragraphs, and with its U(q) repetition supplied as a further conditional reading, in which multiplicity \(q\) of the nonzero spectrum of the classical edge charge is read as \(q\) equal rows, that reading fixes a stacked diagram only when that spectral condition holds, and within the supplied class of Gaussian states the multiplicity is a property of the state, not of the background alone: for q coincident copies of a compact fuzzy factor, with the factor matrices and the cap projector both lifted as tensor products of the identity of the copy space with the one-copy data, and momenta taken copy-invariant, the block-cut edge charge is exactly the tensor product of that identity with the one-copy charge, every nonzero eigenvalue has multiplicity exactly q and the moments \(\operatorname{Tr}Q^{2p}\) acquire the common factor q, verified on explicit small matrices to within roundoff for single caps and for product blocks, so the rule reads the diagram with each row length repeated q times: the stacked diagram of common depth q, to which the exact closed form and the continuum limits of the third of the four preceding paragraphs apply. For generic momenta on the same coincident background the charge has the same rank but its nonzero eigenvalues split apart, so background coincidence alone does not select stacking. For the band-separated caps of the row-length paragraph, whose combined spectrum has multiplicity 1, no stacked diagram is selected: with saddle row lengths, all proportional to \(N\), and depths \((q_1,q_2)\), the stacked moment signature exceeds the multiplicity-1 signature by the fraction \(\sum_f(q_f-1)w_f^p/\sum_f w_f^p\), with \(w_f=k_fm_f(1-m_f)\) and \(k_f\) the edge-length coefficient of the factor, a ratio independent of \(N\), of the saddle parameter and of the row-length proportionality, equal to \(9/17\) at \(p=1\) and \(81/145\) at \(p=2\) for \((q_1,q_2)=(2,1)\), \((m_1,m_2)=(1/2,1/3)\), \(k_f=1\), and equal to \(1\) at equal depths \(2\): the mismatch is of leading order in \(N\) and is therefore not suppressed at the leading semiclassical order at which the rule fixes one dominant diagram. For the product testbed of the first of the four preceding paragraphs with \(q_f\) coincident copies of factor f, the edge-operator spectrum is the one-copy multiset repeated \(q_1q_2\) times irrespective of the state, while the charge multiplicity under copy-invariant momenta is uniform over the whole spectrum, of depth \(q_1q_2\), or of depth \(q_1\) when the momenta are invariant in factor 1 only, so within this class the per-factor stacked diagram \((\ell_1^{q_1},\ell_2^{q_2})\) with \(q_1\neq q_2\) is never the diagram read from the multiplicity of the charge spectrum. These are selection-rule consistency statements for the coincident-copy and band-separated subclasses of the block-cut edge data of the preceding paragraphs, inside the transplanted rule with its single-dominant-diagram hypothesis at leading semiclassical order and its diagonal-cut convention; they leave the two-sided reading of the \(\pm\) pairs untreated, they do not fix the state that sets the multiplicity, and they close nothing in Problem~G.
